\scanpage{5}
\unsure{transforme} hermitiens en hermitiens). On a en \unsure{particulier} un \uline{accouplement} entre \(\Gamma\) et \(C'\), positif etc. On a
\[
\varphi(a)=\lim_{\substack{\deleted{\unsure{x\in C}}\\\deleted{\unsure{1\,\mathrm{fbt}}}\;\inserted{x\xrightarrow{\unsure{\mathrm{fbt}}}1}\\\unsure{\wedge}\\C}}ax
=\lim_{\substack{x\to1\\\unsure{\wedge}\,\unsure{\mathrm{fbt}}\\C}}xa,
\]
(\unsure{limite faible} dans \(C''\)), d'où \deleted{\unsure{g}}
\[
\langle a,f\rangle=\lim_{\substack{x\to1\\\unsure{\wedge}\,\unsure{\mathrm{fbt}}\\C}}\langle ax,f\rangle
\]
\((a\in\Gamma_1,\ f\in C')\), et cet accouplement est compatible avec ordre, involution etc, et prolonge celui donné sur \(C\times C'\).

Supposons maintenant que \(R\) est un anneau de von Neumann,\inserted{et \(C\) dense dans \(R\)} d'où \unsure{canonique} \(\psi:C''\to R\) (\unsure{applique} \(C''\) sur l'adhérence faible de \(C\) dans \(R\)). Évidemment \(\psi(\Gamma)\subset\Gamma_1\). D'autre part,
\[
\varphi\psi\deleted{\unsure{(u)}}=1\quad\text{sur }\Gamma,
\]
car soit \(u\in\Gamma\), \deleted{\illegible{1}} et \(\unsure{v}=\psi(u)\in\Gamma_1\), on a pour \(\forall f\in R_*\)
\[
\langle v,f\rangle=\langle u,f_1\rangle
\]
\((f_1\text{ image de }f\text{ dans }C';\;\text{en fait on a une injection }R_*\subset C')\).
J'affirme que
\[
\langle v,f\rangle=\langle\varphi(\unsure{v}),f_1\rangle
\]
[ce qui prouve \deleted{que} \(\varphi(\unsure{v})\unsure{-}u\) orthogonal à \(R_*\)] cela se voit facilement car
\[
\langle\varphi(\unsure{v}),f_1\rangle=\lim_{x\to1}\langle vx,f\rangle
\]
en prenant ici le filtre des sections \unsure{croissantes des éléments} \(x\unsure{\ll}1\) de \(C\). \unsure{On est ramené} à prouver : un élément de \(\uline{\Gamma}\) orthogonal à \(R_*\) \deleted{\unsure{est nul}} (i.e. dont l'image dans \(R\) est \(0\)) est nul. C'est immédiat

\mspage{6}
